% Teil II: der durchgerechnete Weg.

\subsection*{Gegeben und gesucht}
(a) $f\colon[0,\infty)\to\mathbb R$, $f(x)=\frac{x^2-4}{x^2+1}$;
(b) $f\colon\mathbb R\to\mathbb R$, $f(x)=3+\mathrm e^{-x}$.
Gesucht: Bild $f(D)$, Umkehrfunktion $g$, Grobskizze der Graphen.

\subsection*{a) $f(x)=\frac{x^2-4}{x^2+1}$}
Umschreiben: $f(x)=\frac{x^2+1-5}{x^2+1}=1-\frac5{x^2+1}$. Ableitung
\begin{equation}
  f'(x)=\frac{10x}{(x^2+1)^2}\ge0 \quad (\text{$>0$ für }x>0)
\end{equation}
Also streng wachsend, Bild von $x=0$ bis $x\to\infty$:
$f(0)=1-\frac51=-4$, $f(x)\to1-0=1$. Injektiv, Bild $[-4,1)$.

Stelle nach $x$ um:
\begin{equation}
  y=1-\frac5{x^2+1}\ \Longrightarrow\ \frac5{x^2+1}=1-y\ \Longrightarrow\
  x^2=\frac5{1-y}-1=\frac{4+y}{1-y}.
\end{equation}
Wegen $x\ge0$:
\begin{equation}
  g(y)=\sqrt{\frac{4+y}{1-y}},\qquad y\in[-4,1).
\end{equation}

\subsection*{b) $f(x)=3+\mathrm e^{-x}$}
$\mathrm e^{-x}\in(0,\infty)$, also $f(x)\in(3,\infty)$; $f$ ist streng
monoton fallend (Ableitung $f'(x)=-\mathrm e^{-x}<0$), also injektiv.
Bild $(3,\infty)$.

Stelle um: $y=3+\mathrm e^{-x}\ \Rightarrow\ \mathrm e^{-x}=y-3
\ \Rightarrow\ x=-\ln(y-3)$:
\begin{equation}
  g(y)=-\ln(y-3),\qquad y\in(3,\infty).
\end{equation}

\begin{ergebnis}
\begin{align*}
  \text{a)}&\ f(D)=[-4,1),\quad g(y)=\sqrt{\tfrac{4+y}{1-y}},\ y\in[-4,1);\\
  \text{b)}&\ f(D)=(3,\infty),\quad g(y)=-\ln(y-3),\ y\in(3,\infty).
\end{align*}
\end{ergebnis}

\begin{kontrolle}
a) Probe $f(g(y))=1-\frac5{1+(4+y)/(1-y)}=1-\frac5{(1-y+4+y)/(1-y)}
=1-\frac{5(1-y)}5=y$, und $g(f(x))$ gibt $x$ ($x\ge0$).
b) Probe: $f(-\ln(y-3))=3+\mathrm e^{\ln(y-3)}=y$. Beide rechnen.
\end{kontrolle}

\begin{merksatz}
Bild und Injektivität kommen aus Monotonie; man stellt $y=f(x)$ nach $x$
um und beachtet die Definitionsbedingungen (Wurzel, $\ln$), um die
Umkehrfunktion samt Bildbereich anzugeben.
\end{merksatz}
